\section{Testing Methodology}

\subsection{Overview}
The data pipeline employs a rigorous testing methodology that follows Test-Driven Development (TDD) principles. Each transformation module is thoroughly tested with synthetic data designed to validate specific functionalities and edge cases.

\subsection{Testing Philosophy}
\begin{itemize}
    \item \textbf{Test-First Development:}
    \begin{itemize}
        \item Tests written before implementation
        \item Synthetic data created with known outcomes
        \item Edge cases identified and tested
    \end{itemize}
    
    \item \textbf{Dual Implementation Validation:}
    \begin{itemize}
        \item Parallel testing of Pandas and DuckDB
        \item Identical results required from both
        \item Performance comparison enabled
    \end{itemize}
\end{itemize}

\subsection{Test Suites}

\subsubsection{Meal Data Validation (test\_1\_combine\_meal\_data.py)}
Comprehensive validation of nutritional calculations:

\begin{lstlisting}[style=pythonstyle, caption=Nutritional Calculation Tests]
def test_calculated_values(load_dataframe):
    """Test that calculated nutritional values are correct."""
    df_filtered = load_dataframe.dropna(subset=[
        "total_calories", "Valeur calorique",
        "total_lipids", "Lipides",
        "total_carbs", "Glucides",
        "total_protein", "Protein"
    ])
    
    # Verify calculations with tolerance
    calories_match = np.isclose(
        df_filtered["total_calories"],
        df_filtered["Valeur calorique"] * df_filtered["quantity"],
        rtol=1e-5
    )
    assert calories_match.all(), "Mismatch in calculations!"
\end{lstlisting}

Key test cases:
\begin{itemize}
    \item Nutritional value calculations
    \item Date and time format validation
    \item Unique record ID verification
    \item Missing data handling
\end{itemize}

\subsubsection{Data Aggregation Testing (test\_3\_group\_data.py)}
Comprehensive test data structure:

\begin{lstlisting}[style=pythonstyle, caption=Test Data Setup]
@classmethod
def setUpClass(cls):
    cls.test_dataframe = pd.DataFrame({
        "meal_record_id": [1, 2, 3, 4...],
        "user_id": [1, 1, 2, 2...],
        "date": ["2024-04-01", "2024-04-01"...],
        "Aliment": ["Cacahuètes", "Amandes"...],
        "Valeur calorique": [828, 600...]
    })
\end{lstlisting}

Test scenarios:
\begin{itemize}
    \item Daily aggregations
    \item User-specific totals
    \item Food type grouping
    \item Statistical calculations
\end{itemize}

\subsubsection{Window Function Testing (test\_4\_*\_window\_function\_*.py)}
Three specialized test suites:

\paragraph{Daily Windows (test\_4\_1)}
\begin{itemize}
    \item 4-week rolling averages
    \item Weekday-specific calculations
    \item Date handling validation
\end{itemize}

\begin{lstlisting}[style=pythonstyle, caption=Window Function Test Data]
cls.test_data = pd.DataFrame({
    "date": ["2024-01-04", "2024-01-11"...],
    "total_calories": [2000, 2200...],
    "total_lipids": [20, 22...],
    # Pre-calculated expected results
    "rolling_avg_total_calories": [2000.0, 2100.0...]
})
\end{lstlisting}

\paragraph{User Windows (test\_4\_2)}
\begin{itemize}
    \item User-specific patterns
    \item Time-based aggregations
    \item Rolling calculations
\end{itemize}

\paragraph{Food Type Analysis (test\_4\_3)}
\begin{itemize}
    \item Category-based calculations
    \item Nutritional proportions
    \item Food type patterns
\end{itemize}

\subsection{Test Data Design}

\subsubsection{Synthetic Data Generation}
Carefully crafted test data:

\begin{table}[h]
    \centering
    \begin{tabular}{|l|l|p{8cm}|}
        \hline
        \textbf{Data Type} & \textbf{Example} & \textbf{Purpose} \\
        \hline
        Meals & Cacahuètes, Amandes & Test different food categories \\
        Nutrients & 828 kcal, 20g protein & Validate calculations \\
        Dates & 2024-04-01 & Test temporal patterns \\
        Users & ID: 1,2,3,4 & Test user grouping \\
        \hline
    \end{tabular}
    \caption{Test Data Categories}
\end{table}

\subsubsection{Edge Cases}
Specific scenarios tested:
\begin{itemize}
    \item Missing nutrient values
    \item Zero quantities
    \item Duplicate records
    \item Invalid dates/times
    \item Boundary dates
\end{itemize}

\subsection{Testing Infrastructure}

\subsubsection{Test Base Classes}
Common test infrastructure:
\begin{itemize}
    \item TestDataSetup for shared data
    \item Fixture-based data loading
    \item Cleanup procedures
    \item Error tolerance settings
\end{itemize}

\subsubsection{Assertion Strategies}
Specialized comparisons:
\begin{itemize}
    \item DataFrame equality checks
    \item Floating-point tolerances
    \item Date format validation
    \item Null handling
\end{itemize}

\subsection{Development Workflow}

\subsubsection{TDD Process}
Step-by-step approach:
\begin{enumerate}
    \item Create test data with known outcomes
    \item Write test cases for new functionality
    \item Implement the feature
    \item Verify against both implementations
    \item Document test cases and data
\end{enumerate}

\subsubsection{Quality Gates}
Testing requirements:
\begin{itemize}
    \item All tests must pass
    \item Both implementations must match
    \item Edge cases must be covered
    \item Performance metrics collected
\end{itemize}

This comprehensive testing methodology ensures the reliability and accuracy of the nutritional data pipeline, with particular emphasis on data integrity, calculation accuracy, and implementation consistency between Pandas and DuckDB approaches.
